% !TEX root = ../main.tex

\section{Introduction}
\label{sec:introduction}

Feature selection chooses the features on which a prediction model is trained. A wrapper method judges each candidate subset of features by training the model with it and measuring how well the model predicts. When predictive quality and the number of selected features, called the subset size or \emph{cardinality}, are treated as two separate objectives, the result is not a single subset but a set of trade-offs from which a user can choose~\cite{Xue2016Survey,Jiao2024Survey}. Evolutionary multi-objective algorithms, such as the nondominated sorting genetic algorithm II (NSGA-II)~\cite{Deb2002NSGAII}, are a common way to search for such a set. A run returns a front of subsets in which each larger subset predicts better than every smaller one. Small subsets often matter most in practice: a model with few features is easier to inspect and explain, each feature may be costly to measure or compute, and a deployed model needs only the features it uses. Fronts are therefore often compared at small feature limits, for example with the area under the front (AUF) introduced with the FASTENER method~\cite{Koprivec2020}. For each limit $k$ up to a chosen maximum, such as ten features, AUF takes the best score available with at most $k$ features and averages these scores over the limits.

A score at a feature limit answers two questions at once. The first is whether the method returns any subset within the limit. The second is how well the subsets it returns predict. AUF counts a limit as zero when the front holds no subset with at most that many features, so a zero means that nothing was returned, not that a poor subset was returned. We say that a run \emph{reaches} a limit if its returned front contains a subset within it, and we call the range of subset sizes on the front its \emph{cardinality coverage}. A large AUF difference between two methods is easily read as better search among small subsets, when the method with the lower score may simply not have reached the small sizes. The two readings lead to different conclusions: the first concerns how well the search works among small subsets, the second whether it gets there at all within the available number of model fits.

How a search can fail to reach a small limit is visible in how it starts and how it changes subsets; Section~\ref{sec:metrics} gives the details. Many initializations select each feature with probability one half, so every candidate starts with about half of the features. The opposition-based initialization~\cite{Tizhoosh2005OBL}, as used in the NSGAII-MIIP algorithm~\cite{Li2025}, pairs each such candidate with its opposite, which selects exactly the features the first one does not and so also has about half of them. On a dataset with a few hundred features, such a population almost never contains a subset with ten or fewer features. The search must then move down to small sizes through crossover, mutation, and selection, and a finite budget of model fits may end before it gets there. Once subsets are small, a second effect appears. The balanced mutation of NSGAII-MIIP~\cite{Oh2004,Li2025} keeps the expected subset size unchanged, but when a subset is small relative to the number of features, it changes very few features. Many children then repeat subsets evaluated before. A cache, which stores evaluation results so that a repeated subset needs no new fit, keeps these repeats out of the fit count, so they cost time without appearing in the budget.

The role of subset size at initialization is already recognized. Earlier work allocates initial candidates across size ranges, compares random and informed starts, and adds sparse candidates, that is, subsets with few selected features~\cite{Deniz2019,Kropp2022}. Miyandoab et al.~\cite{Miyandoab2024} first draw a subset size and then select that many features. Cong et al.~\cite{Cong2026} compare sparse and dense starts at a fixed evaluation target and report that they differ in access to compact fronts in a gene-panel application. The unchanged expected size of the balanced mutation is stated in~\cite{Oh2004,Li2025}, and AUF comes from FASTENER~\cite{Koprivec2020}. These studies show that the start affects which subset sizes are available to the search. Our question is how this affects the reading of a reported comparison. How much of a difference at a small feature limit is a difference in reach? When does the lack of reach occur, in terms of the number of features and the budget? What does a sparse start give up among large subsets? And what should be reported so that a reader can tell reach and quality apart?

To answer these questions, we compare three configurations of NSGA-II that share the genetic operators of NSGAII-MIIP: its opposition-based start, in which each candidate begins with about half of the features; a sparse start, in which each candidate has 1 to 30 features; and the complete NSGAII-MIIP, which adds a mutual-information improvement phase to the opposition-based start. We run them on 46 datasets with 16 to 192{,}000 features, with up to three classifiers and five seeds, under a fixed budget of model fits. A pilot on four of these datasets adds larger budgets, starts that cover all subset sizes, and a mutation that always changes a subset. We do not propose a new initialization, mutation, or evaluation measure. We use cardinality-based initialization~\cite{Miyandoab2024}, balanced mutation~\cite{Oh2004,Li2025}, and AUF~\cite{Koprivec2020} as they were introduced, and we examine how initialization and mutation together determine which subset sizes a run returns. The contributions are threefold:
\begin{enumerate}
    \item \textbf{Evidence that a large AUF difference can come almost entirely from reach.}
    We compare the two starts on 46 datasets with 16 to 192{,}000 features and five seeds each. On 39 of these datasets, all three configurations were run with each of three classifiers: a decision tree, a random forest, and LightGBM. With a limit of ten features, the sparse start has a much higher AUF than the opposition-based start. On no dataset does the opposition-based start return a subset within the limit in more runs than the sparse start, and 98 to 100 percent of the AUF difference comes from limits at which the opposition-based start returns no subset at all. Sizes that both starts reach occur only on 1 to 11 datasets per classifier, all with at most 235 features. On them, the mean difference in test F1 lies within $\pm 0.02$ under every classifier. The AUF difference therefore measures mainly whether a run returns any subset with at most ten features, not how well its subsets predict at a given size. The sparse start in turn has a lower best test F1 at any size, by 0.007 to 0.029 on average, so neither start covers the whole range of subset sizes. A pilot shows that starts covering all sizes do not change this: once small subsets are present, selection keeps them and the large subsets disappear.

    \item \textbf{An account of when and why the search stays near its starting sizes.}
    We calculate the probability that an initial population contains a subset with at most a given number of features. For an opposition-based population of 100 candidates and a limit of ten features, it falls below 0.01 from 46 features on. For the balanced mutation, we show that the expected subset size is unchanged while the expected number of changed features is $2k/d$, where $k$ is the current subset size and $d$ the number of features. The runs are consistent with this account. Within the benchmark budgets, the opposition-based search reaches ten features on every dataset with at most 99 features and on none with more than 512; half of its runs reach the limit at about 160 to 230 features. The deficit is one of budget, not a barrier: in the pilot, on four datasets with 176 to 720 features, the opposition-based start reaches the limit in every run when given ten times the budget, after 0.5 to 7.1 times the budget, roughly in proportion to the number of features. The improvement phase of NSGAII-MIIP raises the share of runs that reach the limit by only 0.04 to 0.08. In sparse runs, in which few features change per generation, 96 to 99 percent of the proposed subsets had already been evaluated, and the share of children left unchanged by mutation follows $e^{-2k/d}$. A mutation that always changes a feature makes these searches far cheaper, but not better, because one change to a converged population usually gives a subset evaluated before. The account suggests a mechanism for the observed pattern; it does not establish it as the only cause.

    \item \textbf{Recommendations for reporting comparisons at a feature limit.}
    From these results we derive a short list of quantities to report alongside AUF: whether each run reaches the limit and after how much of the budget, the score at sizes reached by both methods, the best score among larger subsets, the numbers of proposed subsets and new model fits, the elapsed time, and the runs that did not finish.
\end{enumerate}

The results on 8 of the 46 datasets come from an experiment designed for this comparison. The results on the other 38 datasets reuse runs from a broader benchmark of feature-selection methods. We wrote the criteria of the analysis after the results and an exploratory analysis had been read, so we report it as a secondary analysis that describes the tested configurations, budgets, and classifiers, not as a confirmation of a general property of NSGA-II. The pilot was registered before its runs and uses new seeds, but its four datasets had been seen before, so it guides the interpretation and does not confirm it.

The rest of the paper is organized as follows. Section~\ref{sec:related-work} reviews work on initialization, the control of subset size during search, the evaluation of fronts, and repeated candidates, and states what this paper adds. Section~\ref{sec:metrics} defines the problem, AUF, reach, and quality at matched sizes, and analyzes access to small subsets at initialization and how often mutation changes a subset. Section~\ref{sec:design} describes the datasets, the configurations, the statistical analysis, and the pilot. Section~\ref{sec:results} reports the results, Section~\ref{sec:discussion} discusses how they should be read, their limitations, and what to report, and Section~\ref{sec:conclusion} concludes. Appendix~\ref{sec:appendix} examines missing runs and the choice of classifier.

% Provenance of the numbers in this section:
% Pilot: evidence/H14-E02 (fastener-experiments e44fd74).
% H13-E01 (fastener-experiments e023933), copied to evidence/H13-E01/: q1_criterion (contribution 1),
% q2_reach_by_dataset and q2_reach_model (contribution 2), q3_pairs (NSGAII-MIIP), q4_large_end
% (best F1 at any size), q8_stagnation (repeats). Matched-size numbers apply the 50-test-row rule
% (DEC-57, fastener-experiments 354c307): q1_criterion_min_test50 and its _exact variant. Without the
% rule, the forest interval is [-0.044, +0.008]; with it, [+0.0003, +0.0126].
